\documentclass[10pt,a4paper]{article}
\usepackage{amsmath,amssymb,amsthm,graphicx}
\usepackage{amsmath,amssymb,amsthm,hyperref,xfrac}
\providecommand{\eg}{e.g.\ }
\providecommand{\ie}{i.e.\ }
\usepackage[scr=rsfs,cal=euler]{mathalfa}
\multlinegap0pt
\hyphenation{accu-rate applied assump-tion assump-tions decom-pos-ing defined esti-mate esti-mates exten-ded hypothe-ses para-me-ter para-me-ters remain-der required}
\newenvironment{enumeratei}
{\bgroup\def\theenumi{\roman{enumi}}\def\theenumii{\arabic{enumii}}\begin{enumerate}}
{\end{enumerate}\egroup}
\newenvironment{enumeratea}
{\bgroup\def\theenumi{\alph{enumi}}\begin{enumerate}}
{\end{enumerate}\egroup}

\newcommand{\Psfrac}[2]{(\sfrac{#1}{#2})}
\newcommand{\psfrac}[2]{\sfrac{(#1)}{#2}}
\newcommand{\spfrac}[2]{\sfrac{#1}{(#2)}}
\newcommand{\pspfrac}[2]{\sfrac{(#1)}{(#2)}}
\newcommand{\shalf}{\sfrac{1}{2}}

\newcommand{\bnorm}[1]{\bigl\|#1\bigr\|}
\newcommand{\Bnorm}[1]{\Bigl\|#1\Bigr\|}
\newcommand{\babs}[1]{\bigl|#1\bigr|}
\newcommand{\Babs}[1]{\Bigl|#1\Bigr|}

\AtBeginDocument{%
\DeclareMathOperator{\sinn}{sin}\let\sin\sinn
\DeclareMathOperator{\coos}{cos}\let\cos\coos
}

%\usepackage{enumitem}
\usepackage{physics}
%\usepackage[foot]{amsaddr}
\usepackage{bm}
\usepackage{tikz}
\usepackage{mathtools}
\usepackage{scalerel,stackengine}
\stackMath
\usetikzlibrary{arrows.meta}
\usepackage{subcaption}
%\usepackage[style=numeric]{biblatex}
%\addbibresource{main.bib}
%\setlength{\bibitemsep}{0.3\baselineskip}

\newenvironment{mycases}
{\begin{dcases*}}
{\end{dcases*}}
\theoremstyle{plain}
\newtheorem{theorem}{Theorem}[section]
\newtheorem{lemma}[theorem]{Lemma}
\newtheorem{definition}[theorem]{Definition}
\newtheorem{proposition}[theorem]{Proposition}
\newtheorem{corollary}[theorem]{Corollary}
\newtheorem*{corollary*}{Corollary}
\theoremstyle{remark}
\newtheorem{remark}[theorem]{Remark}
\counterwithout{figure}{section}
\counterwithout{equation}{section}

\newcommand{\br}[1]{\left(#1\right)}
\newcommand{\bbr}[1]{\bigl(#1 \bigr)}
\newcommand{\Bbr}[1]{\Bigl(#1 \Bigr)}
\newcommand{\bbbr}[1]{\biggl(#1 \biggr)}
\newcommand{\tpd}{\frac{1}{(2\pi)^{d}}}
\newcommand\pig[1]{\scalerel*[2.35pt]{\big#1}{%
\ensurestackMath{\addstackgap[1.2pt]{\big#1}}}}
\newcommand\pigl[1]{\mathopen{\pig{#1}}}
\newcommand\pigr[1]{\mathclose{\pig{#1}}}
\newcommand{\ak}{\abs{k}}
\newcommand{\ap}{\abs{p}}
\newcommand{\aal}{\abs{\alpha}}
\newcommand{\sbbr}[1]{\bigl[#1\bigr]}
\newcommand{\sBbr}[1]{\Bigl[#1\Bigr]}
\newcommand{\sbbbr}[1]{\biggl[#1\biggr]}
\newcommand{\sbr}[1]{\left[#1\right]}
\newcommand{\half}{\frac{1}{2}}
\newcommand{\p}{\partial}
\newcommand{\jb}[1]{\langle#1\rangle}
\newcommand{\bset}[1]{\bigl\{#1\bigr\}}
\newcommand{\Bset}[1]{\Bigl\{#1\Bigr\}}
\newcommand{\set}[1]{\left\{#1\right\}}
\newcommand{\mytextfrac}[2]{{\textstyle\frac{#1}{#2}}}
\newcommand{\qaq}{\quad\text{and}\quad}
\newcommand{\para}{\parallel}
\newcommand{\ep}{\varepsilon}
\newcommand{\R}{\mathbb{R}}
\newcommand{\C}{\mathbb{C}}
\newcommand{\al}{\alpha}
\newcommand{\intRd}{\int_{\R^d}}
\newcommand{\intRdk}{\int_{\R^d_k}}
\newcommand{\intRdp}{\int_{\R^d_p}}
\newcommand{\lam}{\lambda}
\newcommand{\ceils}[1]{\lceil#1\rceil}
\newcommand{\Qin}{Q_{\mathrm{in}}}
\newcommand{\hin}{h_{\mathrm{in}}}
\newcommand\reallywidehat[1]{%
\savestack{\tmpbox}{\stretchto{%
\scaleto{%
\scalerel*[\widthof{\ensuremath{#1}}]{\kern.pt\mathchar"0362\kern.1pt}%
{\rule{0ex}{\textheight}}%WIDTH-LIMITED CIRCUMFLEX
}{\textheight}%
}{2.4ex}}%
\stackon[-6.9pt]{#1}{\tmpbox}%
}
\title{The argument principle turns a sign condition on the imaginary axis into the uniform Penrose condition}
\author{Marnie Smith}
\date{doi:10.5802/jep.353; CC BY4.0}
\begin{document}
\maketitle
\noindent\textit{Editorial note.} The counted statement is the author's labelled lemma environment at source0.tex lines 1139-1146; its proof environment opens at line 1147, immediately after \textbackslash{}end\{lemma\}, so the pairing is by adjacency and the retained proof range 1147-1189 is the whole environment. The two labelled implications impliesPenrose and impliesPenrose0 used in the proof are declared inside the counted statement itself.\par\medskip
\noindent\textit{Editorial note.} Four environments are retained as context: the uniform Penrose condition (source lines 203-205), which is the conclusion; the directional integral estimate gnorm (source lines 1037-1042); the integral representation Lingoodcoordinates (source lines 1059-1062); and the boundedness and decay lemma (source lines 1092-1106).\par\medskip
\noindent\textit{Editorial note.} This article carries no label-number dump and the retained counter values are left empty, so the printed numbering of the card is the wrapper's default; the identity of the counted item is its author label argumentprinciplelemma.\par\medskip
\noindent\textit{Editorial note.} The retained preamble is the publisher layout adapter, the native font and layout package and the native mathematical macro and theorem block of the pinned source.\par\medskip
\noindent\textit{Editorial scope.} One original result of the pinned publisher article, counted once: the argument-principle reduction of the uniform Penrose condition to a sign condition on the imaginary axis (author label argumentprinciplelemma).\par\medskip
\noindent\textit{Allowed standard prerequisites (editorial).}\begin{itemize}
\item Linearised operator symbol and its low-frequency limit. The symbol is an integral transform of the background, and its limit as the parameter tends to zero is an explicit integral transform evaluated at zero parameter.. They are the two objects on which the sign conditions are imposed.
\item Uniform Penrose condition. The uniform Penrose condition requires that the real part of the symbol stay above a negative bound uniformly in the parameter, the spectral variable and the frequency direction.. It is the conclusion of the statement.
\item Analyticity and the argument principle. The symbol is analytic in the half-plane of the spectral variable, and a curve lying in the image of the imaginary axis that avoids a half-line bounds the argument of the relevant factor.. It is the mechanism that converts the sign condition into the bound.
\end{itemize}\medskip
\begin{align}\label{Penrose}
\inf_{\substack{\Re\lambda \geq 0 \\ k \in \R^d}} \abs{1 + \widetilde{\mathscr{L}}(\lambda, k)} \geq {c_{\mathrm{P}}}
\end{align}

\begin{lemma}\label{gnorm}
Let \( n \geq 0 \), and let \( \alpha \) be a multi-index. Then, for every $k\neq0$,
\begin{align*}
\int_0^\infty\jb{s}^{n}\babs{\p^\alpha_k\widehat g\bbr{\mytextfrac{k}{\abs{k}}s}}\,ds\lesssim_{d,
\delta_0,n,\aal}\norm{g}_{H^{n+{{1/2}}+\delta_0}_{\aal+\ceils{d/2}}}.\end{align*}
\end{lemma}

\begin{align}
\widetilde {\mathscr L}(\lam,k;\hbar)
= \frac{2}{\hbar\ak} \widehat w(k) \int_0^\infty e^{-\lam \sfrac{s}{\ak}} \sin(\mytextfrac{\hbar}{2} \ak s) \widehat g\bbr{\mytextfrac{k}{\ak}s} \, ds. \label{Lingoodcoordinates}
\end{align}

\begin{lemma}\label{boundedness and decay lemma}
Suppose $w\in L^1(\R^d)$ and $g\in H^{3/2+\delta_0}_{\ceils{d/2}}(\R^d)$. Then:\vspace*{-3pt}
\begin{enumeratea}
\item \label{(a)} As $\ak\to\infty$, the quantity $\widetilde {\mathscr L}(\lam,k;\hbar)\to0$, with decay uniform in $\Re\lam\geq0$ and $\hbar\in[0,1]$.

\item \label{(b)}There exists a constant $C_1=C_1(d,\delta_0)$, independent of $\hbar\in [0,1]$, such that for all $\Re\lam\geq0$ and $k\in\R^d$,\vspace*{-3pt}
\begin{align*}
\babs{\widetilde {\mathscr L}(\lam,k;\hbar)}\leq C_1\norm{w}_{L^1}\norm{g}_{H^{3/2+\delta_0}_{\ceils{d/2}}}.
\end{align*}
\item \label{(c)} If in addition $g\in H^{3/2+\delta_0}_{1+\ceils{d/2}}(\R^d)$, then there exists a constant $C_2=C_2(d,\delta_0)$, independent of $\hbar\in [0,1]$, such that for all $\Re\lam\geq0$ with $\lam\neq0$ and $k\in\R^d$,\vspace*{-3pt}
\begin{align*}
\babs{\widetilde {\mathscr L}(\lam,k;\hbar)}\leq C_2\frac{\ak}{\abs{\lam}}\norm{w}_{L^1}\norm{g}_{H^{3/2+\delta_0}_{1+\ceils{d/2}}}.
\end{align*}
\end{enumeratea}
\end{lemma}

\begin{lemma}\label{argumentprinciplelemma}
Suppose \( w \in L^1(\R^d) \), \( g \in H^{3/2+\delta_0}_{1+\lceil d/2 \rceil}(\R^d) \), and that the following implications hold:\begin{align}
\label{impliesPenrose}\forall\hbar\in[0,1],\,\tau\in\R,\,k\in\R^d,\quad \Im\widetilde {\mathscr L}(i\tau,k;\hbar)=0\implies \Re\widetilde {\mathscr L}(i\tau,k;\hbar)>-1,
\end{align}and, for the low-frequency response,
\begin{align}
\label{impliesPenrose0}\forall\tau\in\R,\,\omega\in\mathbb S^{d-1},\quad \Im\widetilde {\mathscr L}_0(i\tau,\omega)=0\implies \Re\widetilde {\mathscr L}_0(i\tau,\omega)>-1.
\end{align} Then the uniform Penrose condition \eqref{Penrose} is satisfied.
\end{lemma}

\begin{proof}
For each $\hbar\in[0,1]$ and $k\in\R^d$, the assumptions on $w$ and $g$ ensure that the map $\lam\mapsto \widetilde {\mathscr L}(\lam,k;\hbar)$ is analytic in the half-plane $\set{\Re\lam>0}$, continuous on $\set{\Re\lam\geq0}$ and, by Lemma~\ref{boundedness and decay lemma}~\eqref{(c)}, tends to zero as $\abs{\lam}\to\infty$. By the argument principle, the number of zeros of the function $\lam\mapsto 1+\widetilde {\mathscr L}(\lam,k;\hbar)$ in $\set{\Re\lam>0}$ is equal to the winding number of the curve\vspace*{-3pt}
\begin{align*}
\Gamma:=\bigl\{\widetilde {\mathscr L}(i\tau,k;\hbar):\tau\in \R\bigr\}
\end{align*}
around the point $-1$. The implication \eqref{impliesPenrose} ensures that $\Gamma$ does not meet the half-line $(-\infty,-1]$. It follows that $1+\widetilde {\mathscr L}(\lam,k;\hbar)$ has no zeros in $\set{\Re\lam\geq0}$. The same argument, applied to $\mu\mapsto1+\widetilde{\mathscr L}_0(\mu,\omega)$ for fixed $\omega\in\mathbb S^{d-1}$ and using \eqref{impliesPenrose0}, shows that $1+\widetilde{\mathscr L}_0(\mu,\omega)$ has no zeros in $\set{\Re\mu\geq0}$.

It remains to show that the infimum in \eqref{Penrose} is positive. For $k\neq0$, $\hbar\in[0,1]$ and $\Re\lam\geq0$, the representation \eqref{Lingoodcoordinates} reads\vspace*{-3pt}
\begin{align*}
\widetilde{\mathscr L}(\lam,k;\hbar)&\hphantom{:}=\widehat w(k)\,\Phi\br{\sfrac{\lam}{\ak},\hbar\ak,\sfrac{k}{\ak}},\\
\Phi(\mu,h,\omega)&:=\int_0^\infty e^{-\mu s}\,\frac{2}{h}\sin\br{\sfrac{hs}{2}}\widehat g(\omega s)\,ds,
\end{align*}
where the factor $2h^{-1}\sin(hs/2)$ is understood as $s$ when $h=0$, so that $\widetilde{\mathscr L}_0(\mu,\omega)=\widehat w(0)\Phi(\mu,0,\omega)$. The function $\Phi$ is continuous on $\set{\Re\mu\geq0}\times[0,\infty)\times\mathbb S^{d-1}$. Indeed, the integrand is continuous and bounded by $s\norm{\widehat g}_{L^\infty}$, so the integral over $[0,R]$ is continuous for every $R>0$, while $\abs{2h^{-1}\sin(hs/2)}\leq s$ and the Cauchy--Schwarz and trace estimates in the proof of Lemma~\ref{gnorm} give, uniformly in $\Re\mu\geq0$, $h\geq0$ and $\omega\in\mathbb S^{d-1}$,\vspace*{-3pt}
\begin{align*}
\int_R^\infty\Babs{e^{-\mu s}\,\frac{2}{h}\sin\br{\sfrac{hs}{2}}\widehat g(\omega s)}\,ds
&\leq\br{\int_R^\infty\hspace*{-1mm}\frac{s^2}{\jb{s}^{3+2\delta_0}}\,ds}^{\sfrac{1}{2}}\br{\int_0^\infty\hspace*{-1mm}\jb{s}^{3+2\delta_0}\abs{\widehat g(\omega s)}^2\,ds}^{\sfrac{1}{2}}\\
&\lesssim_{d,\delta_0} R^{-\delta_0}\norm{g}_{H^{3/2+\delta_0}_{\ceils{d/2}}},
\end{align*}
so $\Phi$ is a uniform limit of continuous functions. Now set\vspace*{-3pt}
\[
F(\mu,r,\omega,\hbar):=1+\widehat w(r\omega)\,\Phi(\mu,\hbar r,\omega),
\]
so that $F(\mu,r,\omega,\hbar)=1+\widetilde{\mathscr L}(r\mu,r\omega;\hbar)$ for $r>0$ and $F(\mu,0,\omega,\hbar)=1+\widetilde{\mathscr L}_0(\mu,\omega)$.

By Lemma~\ref{boundedness and decay lemma}, there exist constants \( K > 0 \) and \( M > 0 \) such that, for \( k\neq0 \) and \( \Re \lam \geq 0 \), if either (i) \( \ak \geq K \), or (ii) \( \abs{\lam} \geq M\ak \), then
\[
\babs{1 + \widetilde{\mathscr L}(\lam,k;\hbar)} \geq \frac{1}{2}
\]uniformly in $\hbar\in[0,1]$, while $\babs{1 + \widetilde{\mathscr L}(\lam,0;\hbar)}=1$. Writing $\lam=r\mu$ and $k=r\omega$, the remaining region corresponds to the compact set \[
\mathcal K:=\left\{
(\mu, r, \omega, \hbar) : \,
\Re\mu\geq0,\,\abs{\mu} \leq M,\,
0\leq r \leq K,\,
\omega\in\mathbb S^{d-1},\,
\hbar \in [0,1]
\right\}.
\]
On this compact set, the function $F$ is continuous, since $\widehat w$ and $\Phi$ are, and it does not vanish, by the first paragraph. Hence
\[\textstyle
{c_{\mathrm{P}}} := \min\set{\sfrac{1}{2},\,\min_{\mathcal K}\abs{F}} > 0
\]
is a lower bound for $\babs{1 + \widetilde{\mathscr L}(\lam,k;\hbar)}$ over $\Re \lam \geq 0$, $k\in\R^d$ and $\hbar \in [0,1]$.
This establishes the uniform Penrose condition \eqref{Penrose}.
\end{proof}

\begin{thebibliography}{99}
\bibitem{r1} J. Bedrossian - “Nonlinear echoes and Landau damping with insufficient regularity”, Tunis. J. Math. 3 (2021) no. 1, p. 121-205
\bibitem{r2} J. Bedrossian, N. Masmoudi \& C. Mouhot - “Landau damping: paraproducts and Gevrey regularity”, Ann. PDE 2 (2016) no. 1, article ID 4, 71 pages
\bibitem{r3} J. Bedrossian, N. Masmoudi \& C. Mouhot - “Landau damping in finite regularity for unconfined systems with screened interactions”, Comm. Pure Appl. Math. 71 (2018) no. 3, p. 537-576
\bibitem{r4} J. Bedrossian, N. Masmoudi \& C. Mouhot - “Linearized wave-damping structure of Vlasov–Poisson in $\mathbb{R}^3$”, SIAM J. Math. Anal. 54 (2022) no. 4, p. 4379-4406
\bibitem{r5} A. Borie, S. Hadama \& J. Sabin - “Scattering for the positive density Hartree equation”, 2025
\bibitem{r6} A. Bove, G. Da Prato \& G. Fano - “An existence proof for the Hartree-Fock time-dependent problem with bounded two-body interaction”, Comm. Math. Phys. 37 (1974), p. 183-191
\bibitem{r7} A. Bove, G. Da Prato \& G. Fano - “On the Hartree-Fock time-dependent problem”, Comm. Math. Phys. 49 (1976) no. 1, p. 25-33
\bibitem{r8} T. Cazenave - Semilinear Schrödinger equations, Courant Lecture Notes in Math., vol. 10, American Mathematical Society, Providence, RI, 2003
\bibitem{r9} J. M. Chadam - “The time-dependent Hartree-Fock equations with Coulomb two-body interaction”, Comm. Math. Phys. 46 (1976) no. 2, p. 99-104
\bibitem{r10} T. Chen, Y. Hong \& N. Pavlović - “Global well-posedness of the NLS system for infinitely many fermions”, Arch. Rational Mech. Anal. 224 (2017) no. 1, p. 91-123
\bibitem{r11} T. Chen, Y. Hong \& N. Pavlović - “On the scattering problem for infinitely many fermions in dimensions $d\ge 3$ at positive temperature”, Ann. Inst. H. Poincaré C Anal. Non Linéaire 35 (2018) no. 2, p. 393-416
\bibitem{r12} J. J. Chong, L. Lafleche \& C. Saffirio - “On the $L^2$ rate of convergence in the limit from the Hartree to the Vlasov-Poisson equation”, J. Éc. polytech. Math. 10 (2023), p. 703-726
\bibitem{r13} C. Collot \& A.-S. de Suzzoni - “Stability of equilibria for a Hartree equation for random fields”, J. Math. Pures Appl. (9) 137 (2020), p. 70-100
\bibitem{r14} C. Collot \& A.-S. de Suzzoni - “Stability of steady states for Hartree and Schrödinger equations for infinitely many particles”, Ann. H. Lebesgue 5 (2022), p. 429-490
\bibitem{r15} E. Grenier, T. T. Nguyen \& I. Rodnianski - “Landau damping for analytic and Gevrey data”, Math. Res. Lett. 28 (2021) no. 6, p. 1679-1702
\bibitem{r16} S. Hadama - “Asymptotic stability of a wide class of stationary solutions for the Hartree and Schrödinger equations for infinitely many particles”, Ann. H. Lebesgue 8 (2025), p. 181-218
\bibitem{r17} A. D. Ionescu, B. Pausader, X. Wang \& K. Widmayer - “Nonlinear Landau damping for the Vlasov-Poisson system in $\mathbb{R}^3$: the Poisson equilibrium”, Ann. PDE 10 (2024) no. 1, article ID 2, 78 pages
\bibitem{r18} L. Landau - “On the vibrations of the electronic plasma”, Acad. Sci. USSR. J. Phys. 10 (1946), p. 25-34
\bibitem{r19} M. Lewin \& J. Sabin - “The Hartree equation for infinitely many particles, II: Dispersion and scattering in 2D”, Anal. PDE 7 (2014) no. 6, p. 1339-1363
\bibitem{r20} M. Lewin \& J. Sabin - “The Hartree equation for infinitely many particles I. Well-posedness theory”, Comm. Math. Phys. 334 (2015) no. 1, p. 117-170
\bibitem{r21} M. Lewin \& J. Sabin - “The Hartree and Vlasov equations at positive density”, Comm. Partial Differential Equations 45 (2020) no. 12, p. 1702-1754
\bibitem{r22} C. Mouhot \& C. Villani - “On Landau damping”, Acta Math. 207 (2011) no. 1, p. 29-201
\bibitem{r23} T. T. Nguyen \& C. You - “Modified scattering for long-range Hartree equations of infinite rank near vacuum”, SIAM J. Math. Anal. 57 (2025) no. 6, p. 6486-6497
\bibitem{r24} T. T. Nguyen \& C. You - “Plasmons for the Hartree equations with Coulomb interaction”, Probab. Math. Phys. 6 (2025) no. 3, p. 913-960
\bibitem{r25} E. M. Stein \& G. Weiss - “Interpolation of operators with change of measures”, Trans. Amer. Math. Soc. 87 (1958), p. 159-172
\bibitem{r26} C. You - “Phase mixing estimates for the nonlinear Hartree equation of infinite rank”, 2024
\end{thebibliography}
\end{document}
